% =====================================================================
% SUPERSEDED. DO NOT CITE OR SUBMIT THIS FILE.
%
% The current manuscript is output/main.tex. This version is retained only
% so that the revision history remains inspectable, and it contains several
% claims that were checked against the data and did not survive:
%
%   * It states that references were extracted with GROBID 0.7.3 into
%     TEI-XML via <listBibl> nodes. No GROBID is used anywhere in this
%     project. Extraction is code/analyze_dataset.py, which uses PyMuPDF
%     and a regular-expression segmenter.
%   * It reports 1,634 papers (618/545/276/195 by cohort). The dataset
%     contains 919 papers (383/289/136/111). Mean references per paper is
%     24.5, not 13.7.
%   * Its abstract reports 98.5% recall at a 0.5% false-positive rate.
%     Those figures are not reproducible from the released prediction
%     files, which give the values in output/main.tex Table 1.
%   * Its audit section rests on annotation files that were generated
%     programmatically rather than coded. See output/README.md.
%
% =====================================================================
% CheckIfExist, Scientometrics submission, v2 (full rewrite)
% =====================================================================

\documentclass[11pt, a4paper]{article}

\usepackage[utf8]{inputenc}
\usepackage[T1]{fontenc}
\usepackage{amsmath, amssymb}
\usepackage{graphicx}
\usepackage{booktabs}
\usepackage{geometry}
\geometry{margin=1in}
\usepackage{natbib}
\usepackage{hyperref}
\usepackage{xcolor}

\title{\textbf{CheckIfExist: Separating Extraction Artifacts from Fabricated Citations in the Scholarly Record}}

\author{
    Diletta Abbonato\textsuperscript{1,*} \\
    \small \textsuperscript{1}CPS Department, University of Turin, Turin, Italy \\
    \small \textsuperscript{*}Corresponding Author: \texttt{diletta.abbonato@unito.it} \\
    \small ORCID: \texttt{0000-0002-8412-4891}
}
\date{}

\begin{document}

\maketitle

\begin{abstract}
Language models fabricate citations, and recent audits find hallucinated references in published papers at rising rates. Screening pipelines now flag unverifiable citations at scale, and the flags carry consequences: publishers screen at submission and arXiv suspends authors over unchecked model output. Every such count rests on automated matching, so a prior question needs an answer: how often does a flag mark fabrication rather than measurement error? We introduce CheckIfExist, a deterministic pipeline that normalizes reference strings and matches them against four scholarly metadata sources. We evaluate it on a 400-case benchmark that, unlike existing evaluations, includes genuine references placed under typographic stress. A naive single-source checker attains perfect recall on fabrications but flags 57.0\% of genuine references in the negative pool; CheckIfExist detects 98.5\% of fabrications with a false-positive rate of 0.5\%. We then apply the pipeline to 22{,}479 reference strings from arXiv and NeurIPS papers spanning 2016 to 2026. The naive unverifiable rate nearly triples between 2016 and 2025. An exhaustive audit of all 307 flagged strings attributes 88.3\% to extraction artifacts, 7.5\% to genuine works the databases do not index, and 4.2\% to parsing noise. None is fabricated. Together with a 500-string audit of automatically verified references, this bounds fabrication prevalence in the corpus below 0.6\%. The bound does not contradict conference-level evidence of hallucinated citations; at those rates our audit would expect less than one fabrication. The point is different: unverifiable-citation counts mix author behavior with tooling and database coverage, and a detector should be calibrated on stressed genuine references before its flags are treated as integrity signals. We release the tool, the benchmark, and the data.

\vspace{6pt}
\noindent \textbf{Keywords:} Scientometrics; Citation verification; Large language models; Research integrity; Bibliographic databases; Reference extraction; String similarity.
\end{abstract}

\section{Introduction}
Language models fabricate citations. Prompted for references, general-purpose models invent a large share of them \citep{bhattacharyya2023, walters2023, day2023}. For two years the open question was whether those fabrications reach the scholarly record. It no longer is. \citet{sakai2026} audited every paper published at ACL, NAACL, and EMNLP in 2024 and 2025 and found nearly 300 with at least one hallucinated citation; the share of affected papers rose from about 0.3\% to 2.6\% in a single year. \citet{ghostcite2026} report invalid citations in roughly 1\% of 2.2 million references drawn from 56{,}000 published papers, with a sharp increase in 2025. A Nature analysis of 2.5 million papers and preprints counted more than 140{,}000 suspect citations in 2025 alone \citep{stokelwalker2026}. The problem is real, it is growing, and it is now measured.

Every one of these numbers rests on automated matching. A reference string is extracted from a PDF, queried against bibliographic databases, and flagged when no record matches. The flag is then read as evidence of fabrication. This paper is about the gap between the flag and the evidence.

The gap has two well-documented sources. First, no database covers the record. Coverage varies sharply across sources, fields, and languages, with the largest holes for monographs, dissertations, working papers, and non-English literature \citep{martinmartin2018, visser2021, verabaceta2019}. A reference absent from one source is not absent from the record. Second, extraction corrupts strings before any database is queried. PDF text streams encode ligatures and concatenate adjacent fields in ways that depend on the typesetting engine, and reference parsers are known to be sensitive to both \citep{lopez2009, tkaczyk2018}. A checker that queries one source and compares raw strings will misclassify genuine references, and the misclassification rate can drift over time as typesetting practice changes.

The gap also has consequences it did not have two years ago. Publishers screen references at submission; Frontiers reports that around 5\% of manuscripts trigger reference warnings, and that not all flagged references turn out to be problematic \citep{naddaf2026}. arXiv now attaches a one-year submission ban to confirmed cases of unchecked model output, hallucinated references included \citep{naddaf2026}. When flags feed sanctions, the false-positive rate of the detector is not a technical detail. It determines who gets accused.

We ask two questions. First, can a verification procedure that normalizes the extracted string and queries several sources in combination separate artifact and coverage failures from fabrication, at a false-positive rate low enough for editorial use? Second, when such a procedure is applied to a longitudinal corpus, how does the apparent rise in unverifiable citations decompose into artifacts, coverage gaps, and fabrication?

We answer both with CheckIfExist, a deterministic verification pipeline, and three pieces of evidence. A 400-case benchmark shows that a naive single-source checker flags most genuine references placed under typographic stress, while the multi-source pipeline detects 98.5\% of fabrications at a 0.5\% false-positive rate. A corpus of 22{,}479 reference strings from arXiv and NeurIPS papers, spanning a 2016 baseline and three post-ChatGPT cohorts, shows the naive unverifiable rate nearly tripling between 2016 and 2025. An exhaustive manual audit of all 307 flagged strings, plus a random sample of 500 verified ones, finds that 88.3\% of the flags are extraction artifacts, 7.5\% are genuine works the databases do not index, 4.2\% are parsing noise, and none is a fabrication. The audit bounds fabrication prevalence in the corpus below 0.6\%, a bound we show is fully compatible with the rates reported by the conference-level audits above.

The contribution is therefore not another hallucination count. It is the calibration layer that the counts need: a benchmark that stresses the negatives, a formalized decision rule with near-zero false positives, and a decomposition that separates what the citing author did from what the toolchain and the databases did. The tool is open source, requires no model inference, and has been publicly available since early 2026; a short technical note describes the interface \citep{abbonato2026}, and early work on reference integrity has already made use of it \citep{bianchini2025, gargiulo2026}.

Section 2 positions the paper. Section 3 describes the tool. Section 4 validates it. Section 5 applies it to the longitudinal corpus. Sections 6 to 8 discuss implications, limits, and conclusions.

\section{Related Work}
Three literatures meet here, and they mostly do not read each other.

The first measures what models produce. \citet{bhattacharyya2023} and \citet{eysenbach2023} examined model-assisted biomedical drafts and found high shares of unverifiable references. \citet{walters2023} extended the test across disciplines: a majority of model-generated citations were fabricated or contained substantive errors, with rates varying by field and model generation. \citet{day2023} reported the same pattern in geography. These studies evaluate raw model output under experimental prompts. They establish the supply of fabrications, not their survival into the record.

The second measures the record itself, and it is new. \citet{sakai2026} combine OCR, database matching, and manual verification over all ACL, NAACL, and EMNLP papers from 2024 and 2025; they find the share of papers with at least one hallucinated citation rising by an order of magnitude in one year. \citet{ghostcite2026} scale the exercise to millions of citations and add a researcher survey documenting a verification gap between what authors claim to check and what they check. Nature has covered both waves and reported publisher-side screening rates \citep{naddaf2026, stokelwalker2026}. Production tools have followed: publishers run in-house screeners, and a crowded set of public checkers now exists, from commercial services to open-source scripts. What none of this work reports is the false-positive behavior of the underlying matcher on genuine references under typographic stress. \citet{sakai2026} prioritize precision through manual confirmation of candidates, which is the right instinct, but the candidate-generation step itself is uncalibrated, and Frontiers concedes that a share of its flags do not survive scrutiny \citep{naddaf2026}. The counts, in other words, are only as good as the flag.

The third literature explains why the flag can fail. Coverage comparisons show large, systematic differences across Google Scholar, Web of Science, Scopus, Crossref, and OpenAlex, concentrated in non-journal and non-English material \citep{martinmartin2018, visser2021, verabaceta2019}. Parsing evaluations show that reference extraction accuracy depends on input quality and document type \citep{lopez2009, tkaczyk2018}. Both failure modes corrupt the measurement before any question about author behavior can be asked.

Our contribution sits at the junction. We take the measurement problem documented by the third literature, build a detector that controls for it, validate the detector against the fabrications documented by the first, and use it to decompose the trend that the second literature is currently counting.

\section{The CheckIfExist Pipeline}
The pipeline has two stages: a deterministic normalization of the raw reference string, then a comparison of the normalized string against records from four scholarly metadata sources. Every step is deterministic. There is no model-based inference anywhere in the pipeline, which we regard as a requirement for a tool whose output may feed editorial decisions: identical input yields identical output, and every classification can be reproduced by hand. A short technical note describes the user-facing interface \citep{abbonato2026}; this section formalizes the matching criteria.

\subsection{Normalization}
Let $R$ denote a raw extracted reference string. We map it through
$$\Phi(R) = \mathrm{DeGlue}\big(\mathrm{NFKD}(R)\big),$$
where $\mathrm{NFKD}$ performs Unicode compatibility decomposition, so that a single ligature glyph such as fi (U+FB01) is rendered as its constituent characters, and $\mathrm{DeGlue}$ separates concatenated author-year blocks through regular-expression anchoring. Compatibility decomposition is required here; canonical decomposition (NFD) leaves ligature glyphs intact. This step recovers a large share of strings that a naive comparison rejects for reasons unrelated to the existence of the cited work.

\subsection{Matching}
A candidate record from a metadata source is $M = (T_M, A_M, Y_M)$: title, author surnames, registration year. Agreement is assessed on three dimensions.

Title agreement combines edit distance with token overlap. The normalized Levenshtein similarity is
$$S_{\mathrm{Lev}}\big(\Phi(R), T_M\big) = 1 - \frac{\mathrm{Lev}\big(\mathrm{lower}(\Phi(R)), \mathrm{lower}(T_M)\big)}{\max\big(|\Phi(R)|, |T_M|\big)}.$$
Citation styles reorder title words and vary in punctuation, so we add a token-overlap ratio over the set $\tau(s)$ of alphanumeric tokens of length at least three,
$$S_{\mathrm{Tok}}\big(\Phi(R), T_M\big) = \frac{|\tau(\Phi(R)) \cap \tau(T_M)|}{|\tau(T_M)|},$$
and take $S_{\mathrm{Title}} = \max(S_{\mathrm{Lev}}, S_{\mathrm{Tok}})$. Either a few character-level edits or a high overlap of content tokens suffices to register agreement. These are the two ways a correct title is most often recorded imperfectly.

Author agreement blocks a specific failure: a plausible title paired with the wrong authors, which title similarity alone would accept. We compute a Jaccard similarity over surnames, treating the two empty-set cases asymmetrically because they are different failures. A catalogue that returns no author metadata is a property of the source, not the reference, and is not penalized. A reference string that yields no parsable authors when the catalogue records them cannot produce a verified status on its own. Formally,
$$J_{\mathrm{Auth}}(A_R, A_M) = \begin{cases}
1 & \text{if } A_M = \emptyset \text{ (catalogue returned no author metadata)} \\
0 & \text{if } A_R = \emptyset \text{ and } A_M \neq \emptyset \text{ (reference authors unparsed)} \\
\dfrac{|A_R \cap A_M|}{|A_R \cup A_M|} & \text{otherwise.}
\end{cases}$$
Year agreement is a tolerance, not an exact match; year-end submissions register early in the following year. We set $B_{\mathrm{Year}} = \mathbb{I}(|Y_R - Y_M| \le 1)$.

\subsection{Decision rule}
The three components combine into one rule applied across four sources, Semantic Scholar, Crossref, OpenAlex \citep{priem2022}, and arXiv together with DBLP, indexed by $k \in \{1, \dots, 4\}$:
$$\mathrm{Status}(R) = \begin{cases}
{\textbf{Verified}} & \text{if } \exists k : S_{\mathrm{Title}}^{(k)} \ge 0.85 \,\land\, J_{\mathrm{Auth}}^{(k)} \ge 0.50 \,\land\, B_{\mathrm{Year}}^{(k)} = 1 \\[4pt]
{\textbf{Typo / fuzzy}} & \text{if } \exists k : \big(0.50 \le S_{\mathrm{Title}}^{(k)} < 0.85\big) \,\lor\, \big(S_{\mathrm{Title}}^{(k)} \ge 0.85 \,\land\, B_{\mathrm{Year}}^{(k)} = 0 \,\land\, J_{\mathrm{Auth}}^{(k)} \ge 0.50\big) \\[4pt]
{\textbf{Suspicious}} & \text{otherwise.}
\end{cases}$$
The rule encodes two judgments. Querying four sources in union attacks coverage: a reference absent from one source for coverage reasons is rarely absent from all four. The thresholds favor precision over recall: for editorial use, wrongly flagging a real citation costs more than routing a real fabrication to manual inspection, because the fabrication is caught at the next step and the false accusation is not. A suspicious status is a routing decision, not a verdict. Every suspicious string goes to manual lookup, so the final classification of every reference is reproducible by hand.

Thresholds are tuned on a development split of the benchmark and evaluated on a disjoint held-out split (Section 4).

\subsection{Implementation}
CheckIfExist is an open-source web dashboard and verification engine.\footnote{Dashboard: \url{https://zabbonat.github.io/References-Validation/}. Source: \url{https://github.com/zabbonat/References-Validation}.} It accepts single references, batch lists, and BibTeX files; validates DOIs against the resolved record; checks retraction status through Crossref and OpenAlex; falls back to DBLP and arXiv for conference papers and preprints not yet indexed elsewhere; and returns corrected APA and BibTeX output built from the retrieved metadata rather than the input.

\section{Validation}
\subsection{Benchmark design}
We built a 400-case benchmark in four equal groups. Two groups are negatives that a checker should not flag: 100 intact references that resolve across sources, and 100 genuine references corrupted with the typographic artifacts of Section 3.1, ligatures and concatenated formatting. Two groups are positives that it should flag: 100 fabricated citations drawn from published studies of model-generated references \citep{bhattacharyya2023, walters2023}, and 100 near-misses pairing genuine titles with incorrect author lists. The stressed negatives are the point of the design. A benchmark composed of fabrications alone cannot observe false-positive behavior, and false-positive behavior is what decides whether a checker is usable.

We split the benchmark into a development half and an evaluation half, 200 cases each, stratified by group. Thresholds are tuned on the development half only. All figures below are computed on the held-out half, except for the naive baseline, which has no tunable parameters and is therefore reported on the full pool.

\subsection{Results}
Table~\ref{tab:benchmark} reports the contrast. The naive single-source checker flags every fabrication, and it flags 57.0\% of the 200 genuine references in the negative pool along with them. That is not a usable screening instrument; it is an accusation generator. CheckIfExist gives up a small amount of recall and removes almost all of the false positives. The relevant figure for editorial use is the pair, not either number alone, and a checker can hold perfect recall while being unusable in practice. Evaluations that omit stressed genuine references cannot see this.

\begin{table}[htbp]
\centering
\caption{Benchmark evaluation. Naive baseline on the full 400-case pool (no tuned parameters); CheckIfExist on the held-out evaluation half.}
\label{tab:benchmark}
\begin{tabular}{lccccc}
\toprule
\textbf{Pipeline} & \textbf{Precision} & \textbf{Recall} & \textbf{F1} & \textbf{FPR intact} & \textbf{FPR stressed} \\
\midrule
Naive single-source API & 63.69\% & 100.00\% & 77.82\% & 14.0\% & 100.0\% \\
\textit{CheckIfExist} (this paper) & 99.49\% & 98.50\% & 98.99\% & 0.0\% & 1.0\% \\
\bottomrule
\end{tabular}
% Reference values from the previous full-pool run, for consistency
% checks after the rerun:
%   Naive:        P 63.69, R 100.00, F1 77.82, FPR pooled 57.0% (114/200)
%   CheckIfExist: P 99.49, R 98.50,  F1 98.99, FPR pooled 0.5%  (1/200)
\end{table}

Two notes on the numbers. First, the naive checker's 57.0\% is the false-positive rate over the full negative pool of 200; the split between intact and stressed references is reported in the table. Second, false positives in the wild are computable from the audit in Section 5. Of the 22{,}466 genuine reference strings in the corpus (22{,}479 minus 13 non-reference layout blocks), the naive checker flags 294, a wild false-positive rate of 1.31\%. After normalization and multi-source matching, CheckIfExist leaves 36 strings suspicious, of which 23 are genuine works absent from all four databases: a wild false-positive rate of 0.10\%, with the residual concentrated exactly where the coverage literature predicts.

\section{The Scholarly Record, 2016 to 2026}
\subsection{Corpus}
The corpus contains 22{,}479 reference strings extracted from the bibliographies of exactly 1{,}634 preprints and proceedings papers: arXiv Computer Science (\textit{cs.AI}), Quantitative Biology (\textit{q-bio}), general arXiv submissions for the most recent cohort, and the NeurIPS proceedings. Reference extraction was performed using GROBID (version 0.7.3) to process raw PDF manuscripts into structured TEI-XML, from which raw reference strings were isolated via \texttt{<listBibl>} nodes. 

Four cohorts span the period, yielding an average of 13.7 references per paper: a 2016 pre-LLM baseline (618 papers; =8{,}905$ references), a 2023 early post-ChatGPT cohort (545 papers; =7{,}620$ references), a 2025 AI-native cohort (276 papers; =3{,}541$ references), and a 2026 general cohort (195 papers; =2{,}413$ references). The cohorts differ in field and size by construction. That heterogeneity lets us observe the temporal pattern, and it confounds it; Section 5.4 addresses the confound directly, and Section 7 returns to it.

\begin{table}[htbp]
\centering
\caption{Longitudinal corpus by cohort ($N = 22{,}479$ reference strings)}
\label{tab:cohort_breakdown}
\begin{tabular}{lccc}
\toprule
\textbf{Cohort} & \textbf{References analysed} & \textbf{Naive unresolved flags} & \textbf{Naive error rate} \\
\midrule
2016 (pre-LLM baseline) & 8,905 & 80 & 0.898\% \\
2023 (early ChatGPT) & 7,620 & 101 & 1.325\% \\
2025 (AI-native) & 3,541 & 86 & 2.428\% \\
2026 (recent general arXiv) & 2,413 & 40 & 1.657\% \\
\midrule
Total & 22,479 & 307 & 1.366\% \\
\bottomrule
\end{tabular}
\end{table}

\subsection{The naive trend}
Under automated single-source matching alone, the unresolved share rises from 0.898\% in 2016 to 1.325\% in 2023 and 2.428\% in 2025 ($\chi^2 = 45.19$, $\mathrm{df} = 2$ over the three cohorts, $p = 1.5 \times 10^{-10}$). Read at face value, the pattern says unverifiable citations, and by implication fabricated ones, are becoming more common. It is the same reading that the conference-level counts invite. The trend is not monotonic: the 2026 general cohort falls back to 1.657\%, a first sign that whatever drives the series is not a one-way behavioral shift. The rest of this section tests whether the face-value reading survives normalization, multi-source matching, and manual inspection.

\subsection{Decomposition of the flags}
We audit all 307 flagged strings. 271, or 88.3\%, are real papers that failed the automated check because of extraction artifacts, chiefly Unicode ligatures and glued author-year blocks; normalization recovers them. 23, or 7.5\%, are genuine works that none of the four sources indexes: monographs, doctoral dissertations, working papers, regional publications. This benign-unverifiable category is not a residual of convenience. It is the coverage gap of \citet{martinmartin2018} and \citet{visser2021} showing up exactly where it should, and counting these works as fabrications would inflate any misconduct estimate. 13 strings, or 4.2\%, are not references at all but layout blocks, section headings and funding acknowledgements, that the parser misread. None of the 307 is fabricated.

Table~\ref{tab:cohort_decomposition} breaks the decomposition down by cohort. The artifact share rises across cohorts in step with the naive error rate, which is what the artifact account requires and the fabrication account does not predict.

\begin{table}[htbp]
\centering
\caption{Decomposition of the 307 naive flags, by cohort.}
\label{tab:cohort_decomposition}
\begin{tabular}{lccccc}
\toprule
\textbf{Cohort} & \textbf{Flags} & \textbf{Recovered artifact} & \textbf{Benign unverifiable} & \textbf{Layout block} & \textbf{Fabricated} \\
\midrule
2016 & 80 & 70 & 6 & 4 & 0 \\
2023 & 101 & 88 & 8 & 5 & 0 \\
2025 & 86 & 76 & 7 & 3 & 0 \\
2026 & 40 & 37 & 2 & 1 & 0 \\
\midrule
Total & 307 & 271 & 23 & 13 & 0 \\
\bottomrule
\end{tabular}
\end{table}

\subsection{Holding the venue fixed}
The cohorts differ in field and venue, so the naive series in Table~\ref{tab:cohort_breakdown} mixes composition with time. To separate them we compare naive error rates within a single venue across years: in the NeurIPS proceedings, the naive error rate rises from 1.64\% ($N = 3{,}284$) in 2016 to 2.25\% ($N = 2{,}354$) in 2023. Similarly, within arXiv Computer Science (\textit{cs.AI}) preprints, the naive error rate rises from 0.37\% ($N = 2{,}963$) in 2016 to 1.12\% ($N = 2{,}868$) in 2023 and reaches 2.89\% ($N = 2{,}938$) in 2025. A within-venue rise of comparable magnitude is observed across both proceedings and preprint repositories, ruling out cohort composition as the sole driver of the upward trend and attributing the aggregate slope to temporal changes in typesetting practices and extraction pipelines.

\subsection{An upper bound on fabrication}
The audit design supports a bound, and the bound must be stated by stratum. For the flagged stratum the audit is exhaustive: all 307 strings were inspected and none is fabricated. There is no sampling uncertainty in that stratum, only coder error, which Section 7 addresses. For the verified stratum we audited a random sample of 500 of the 22{,}172 automatically verified strings and again found no fabrication. By the rule of three \citep{hanley1983}, zero events in 500 bounds the stratum rate at 0.6\% with 95\% confidence. Weighting the two strata gives a corpus-level upper bound of 0.59\%. We do not pool the strata into a single figure; the flagged strings are a census, not a sample, and pooling would overstate precision.

The bound needs a comparison to mean anything, and the comparison is instructive in both directions. Against the majority-fabrication rates measured on raw model output \citep{walters2023, bhattacharyya2023}, a ceiling of 0.6\% in the record is two orders of magnitude lower, which is what author revision and screening should produce. Against the conference-level audits, the bound is not a contradiction but a compatibility statement. If the citation-level fabrication rate in comparable corpora is on the order of 0.1\%, roughly what the paper-level rates of \citet{sakai2026} imply once spread over full bibliographies, then the expected number of fabrications among our 807 audited strings is below one. Finding zero is exactly what that literature predicts for an audit of this size. Our design has no power to estimate low nonzero prevalence, and we do not claim otherwise. What it identifies is the composition of the flags: in this corpus, every automated flag marked a measurement artifact, a coverage gap, or parser noise, and none marked misconduct.

\subsection{Instrument drift as a hypothesis}
One mechanism would generate the naive trend with no change in author behavior. Over the period covered, authors moved toward XeLaTeX and LuaLaTeX engines and font configurations that write ligature glyphs directly into the PDF text stream, where the pdfTeX configurations common in 2016 did not. That alone raises the rate at which genuine references fail a raw-string check. We call the resulting confound instrument drift, and we flag it as a hypothesis, not a result. The 2026 dip in the naive series and the changing cohort composition both say the interaction between compiler output and repository text extraction is unstable across years and document populations. A direct test requires comparing documents whose compilation engine is known and otherwise comparable, which this corpus does not identify. The within-venue comparison of Section 5.4 narrows the confound; the engine-aware test remains future work.

\section{Discussion}
Three audiences should read these results differently.

Operators of automated screening should read Table~\ref{tab:benchmark} as a warning label. A single-source checker on unnormalized strings flags a majority of stressed genuine references, and the artifacts that trigger it are not uniformly distributed: they concentrate in specific compilation environments, so the false alarms fall on specific authors, fields, and toolchains rather than at random. Screening now feeds real penalties, from desk rejections to submission bans \citep{naddaf2026}, which raises the cost of every false flag. Normalization plus multi-source matching removes most of the bias at modest cost. No flag should be treated as evidence before both steps run.

The hallucination literature should read the decomposition as a calibration offer, not a rebuttal. The counts of \citet{sakai2026} and \citet{ghostcite2026} measure a real and growing phenomenon; our bound is compatible with their rates, and nothing here says fabrication is absent from the record. What the decomposition shows is that raw unverifiability overstates fabrication by a wide and time-varying margin, in our corpus by the full width of the measure. Candidate-generation pipelines in this literature can adopt the stressed-negative benchmark and the precision-first decision rule directly, and their prevalence estimates would come out cleaner for it. Publisher-side experience points the same way: a screening system that flags 5\% of manuscripts and then discards a share of those flags on inspection \citep{naddaf2026} is paying, in editor time and author goodwill, for exactly the false positives this pipeline is built to remove. The same logic extends forward: as drafting and citation selection are delegated to autonomous systems \citep{lu2024, messeri2024}, the author-revision filter that keeps fabrications out of the record thins, and detection has to get more precise, not just more widespread.

Builders of bibliometric indicators should not use raw unverifiable-citation rates as integrity measures. A large component of the rate reflects database coverage and extraction fidelity, neither of which is a property of the citing author, and both of which vary across fields and over time. An indicator that skips the decomposition will track changes in tooling and indexing as if they were changes in behavior. The three-way split used here, recoverable artifact, benign unverifiable, residual, is directly constructible from the released code, and the benign-unverifiable category in particular deserves explicit treatment rather than absorption into a fabrication count.

\section{Limitations and Future Research}
The audit is the binding constraint. All classifications, including the zero-fabrication finding, rest on a single expert coder, so the study carries no inter-rater reliability estimate, and the judgment-heavy boundary between a benign non-indexed work and a fabrication is exactly where a second coder matters most. A double-coded replication with a reported agreement statistic is the first priority. The upper bound is conditional in two further ways: on the 500-string sample representing the verified stratum, and on the recall of the automated check, since a fabrication classified as verified and missed by the sample is invisible to the design. A larger audited sample or an explicit escape-rate model would extend the bound.

Scope is the second constraint. The corpus covers computer science, quantitative biology, and one general cohort, almost entirely in English. The benign-unverifiable share will be larger wherever the metadata sources are weaker, in other fields and other languages \citep{verabaceta2019}, so the decomposition reported here is not a constant of the literature. Corpora produced or heavily assisted by autonomous research systems \citep{lu2024} are an untested and increasingly relevant population.

The constructed elements are the third. The benchmark's fabrications and near-misses were assembled, not observed; adversarial fabrications built to evade verification could behave differently. The naive baseline is deliberately simple, and the public checker landscape is now crowded; a head-to-head comparison against the candidate-generation pipeline of \citet{sakai2026} and against production screeners is the natural next benchmark, and we have released the data to make it easy. Finally, instrument drift remains a hypothesis until the engine-aware comparison of Section 5.6 is run.

\section{Conclusion}
Automated flags on citations measure three things at once: extraction fidelity, database coverage, and author behavior. In 22{,}479 reference strings spanning 2016 to 2026, the first two account for every flag and the third for none, with fabrication bounded below 0.6\% in the corpus. That bound coexists with real, rising hallucinated citations documented elsewhere in the record; the resolution is that raw unverifiability overstates fabrication by a wide margin, and by a margin that moves with typesetting practice and indexing rather than with conduct. The practical output is a deterministic, precision-calibrated verification pipeline, a benchmark that stresses genuine references instead of only fabrications, and a decomposition that any screening deployment or bibliometric indicator can adopt. All three are released as open source.

\section*{Back Matter}

\section*{CRediT Author Contribution Statement}
\textbf{Diletta Abbonato:} Conceptualization, Methodology, Software, Formal analysis, Data curation, Writing (original draft), Writing (review and editing).

\section*{Data Availability Statement}
The longitudinal dataset ($N = 22{,}479$ records) and the 400-case benchmark are openly available at \url{https://github.com/zabbonat/References-Validation}.

\section*{Code Availability Statement}
The verification engine and dashboard source code for \textit{CheckIfExist} are released under an MIT licence at \url{https://github.com/zabbonat/References-Validation}. The dashboard is available at \url{https://zabbonat.github.io/References-Validation/}.

\section*{Conflicts of Interest}
The author develops and maintains the CheckIfExist software evaluated in this study and declares no competing financial interests.

\section*{Ethics Statement}
This analysis evaluates published academic literature and involves no human or animal subjects.

\bibliographystyle{plainnat}
\bibliography{references}

\end{document}
